\documentclass[11pt]{article}
\usepackage{amsmath, amssymb, amsthm, geometry, bm, booktabs}
\usepackage[hidelinks]{hyperref}
\geometry{margin=1in}

\title{Pair Correlation and Spectral Correspondence in Primorial Lattices:\\
Toward a Geometric Bridge Between Discrete Gap Statistics and the Riemann Zeta Zeros}
\author{Rob Miller}
\date{\today}

% RAPF Paper 5 — the spectral bridge paper
% Framework by Rob Miller. LaTeX scaffold by Rebecca/Hermes.

\newtheorem{theorem}{Theorem}[section]
\newtheorem{lemma}[theorem]{Lemma}
\newtheorem{corollary}[theorem]{Corollary}
\newtheorem{proposition}[theorem]{Proposition}
\newtheorem{conjecture}[theorem]{Conjecture}
\newtheorem{definition}[theorem]{Definition}
\theoremstyle{remark}
\newtheorem{remark}[theorem]{Remark}

\begin{document}
\maketitle

\begin{abstract}
Building on the unified structural framework established in Papers~1--4,
this paper initiates the spectral bridge between the discrete information
geometry of primorial lattices and the analytic theory of the Riemann
zeta function. We pursue three parallel objectives:
\begin{enumerate}
    \item \textbf{Empirical pair correlation:} Extract normalized gap
    statistics from the twin and triplet census datasets (Papers~2 and~3)
    and compare against the Montgomery--Bogomolny--Keating GUE pair
    correlation model.
    \item \textbf{Explicit formula bounds:} Relate short-interval prime
    $k$-tuple counts in primorial ranges to sums over zeta zeros via the
    von Mangoldt explicit formula, using the discrete Taylor penalties
    from Paper~4 to bound the error contributions.
    \item \textbf{RG fixed-point analysis:} Strengthen the heuristic
    evidence for the $1/\sinh^2$ fixed-point conjecture with leading-order
    error bounds under the logarithmic primorial flow.
\end{enumerate}
This paper represents a shift from phenomenological framework-building
to testable, predictive spectral analysis. The computational results
provide new empirical evidence that the primorial lattice repulsion
exhibits the same structural uniformity found in the zero-spacing
statistics of the Riemann zeta function.
\end{abstract}

\noindent \textbf{MSC 2020:} 11N05, 11M26, 58J50

%================================================================================
\section{Introduction}
%================================================================================
% TODO:
% - Recap of Papers 1-4: discrete Fisher metric, census, non-collapse, Taylor
% - The gap in the architecture: from finite-scale structure to zeta zeros
% - Montgomery's pair correlation (1973) and the GUE connection
% - What this paper does: empirical pair-correlation + explicit formula bounds + RG bounds
% - Scope: this is a bridge paper, not a proof of RH or TPC

%================================================================================
\section{Pair Correlation on Primorial Lattices}
%================================================================================

\subsection{Normalization of Gap Data}
% TODO: Normalizing gap sequences for comparison with GUE/RMT predictions

\subsection{Empirical Pair-Correlation Function}
% TODO: Streaming histogram algorithm for n=8 twins / n=7 triplets
% TODO: Plot against 1 - (sin(pi u) / (pi u))^2

\subsection{Comparison to Montgomery's GUE Model}
% TODO: Quantitative comparison — KS distance, chi-squared, etc.

%================================================================================
\section{The Explicit Formula Bridge}
%================================================================================

\subsection{The Riemann--von Mangoldt Explicit Formula}
% TODO: Standard explicit formula relating psi(x) to zeros

\subsection{Bounding Error Terms via Taylor Penalties}
% TODO: How the discrete penalties from Paper 4 bound the zero contributions

\subsection{Implications for Zero-Free Regions}
% TODO: Connection between non-collapse and controlled zero-free regions

%================================================================================
\section{Leading-Order Bounds on the RG Fixed-Point}
%================================================================================

\subsection{Partial Analysis of the O(log p_{n+1}/log P_n) Error}
% TODO: Rigorous leading-order error bounds

\subsection{Trapping the Deviations}
% TODO: Geometric dampening field for Taylor penalty deviations

%================================================================================
\section{Predictive Statistics at Higher Scales}
%================================================================================

\subsection{Maier-Modulated PMF Fit}
% TODO: Fitting M(k) modulation function to census data

\subsection{Prediction for n=9 Twin CV}
% TODO: Hard-coded prediction with error bars

%================================================================================
\section{Discussion}
%================================================================================
% TODO: What the empirical and theoretical results mean together
% TODO: Honest assessment of what is proven vs. suggested

%================================================================================
\section{Conclusion}
%================================================================================
% TODO: Summary of findings, open problems, next steps

%================================================================================
\begin{thebibliography}{99}

\bibitem{montgomery}
H. L. Montgomery, \emph{The pair correlation of zeros of the zeta function},
Proc. Sympos. Pure Math. \textbf{24} (1973), 181--193.

\bibitem{BK96}
E. B. Bogomolny and J. P. Keating, \emph{Gutzwiller's trace formula and the
statistical properties of the Riemann zeta function}, Nonlinearity \textbf{9}
(1996), 911--935.

\bibitem{paper1}
R. Miller, \emph{From Continuous Geometry to Discrete Lattices: An Empirical
Anatomy of Failed and Corrected Information-Geometric Models of Twin Primes},
Paper 1, 2026.

\bibitem{paper2}
R. Miller, \emph{Equidistribution of Twin Primes Across Admissible Residue
Classes: A Computational Census from $n=3$ through $n=8$}, Paper 2, 2026.

\bibitem{paper3}
R. Miller, \emph{Discrete Information Geometry of Prime Triplets: The
Modulo-30 Arithmetic Lattice and Variance Suppression}, Paper 3, 2026.

\bibitem{paper4}
R. Miller, \emph{A Unified Structural Framework of Prime Constellations:
Induction, Renormalization Flow, and Information Non-Collapse}, Paper 4, 2026.

\end{thebibliography}

\end{document}
